\section{From finite observations to a low-velocity asymptotic state}
\label{sec:low-velocity-tail}

A stable missing-velocity inverse can be composed with
a stable evolution toward the singular end in a
uniformly subcritical, small-data region. The following
estimate gives one explicit overlap between those
two restrictions. The asymptotic expansion itself
belongs to the established low-velocity Gowdy
theory \cite{ringstrom,ringstromavtd,libounces}.
Here the estimate controls differences between
limiting states by differences between finite
measurements.

Let $r=\tau-\tau_0\ge0$, fix $0<v_*<1$, and set
\[
 V=P_\tau,\quad W=e^{-\tau}P_\theta,\quad
 X=e^P Q_\tau,\quad Y=e^{P-\tau}Q_\theta,\qquad
 Z=(V-v_*,W,X,Y).
\]
Use the Sobolev norms with weights
$\langle k\rangle^s=(1+k^2)^{s/2}$ and normalized
Fourier coefficients on the $2\pi$ circle.
Let $A_s=2^s(1+2(1+1/(2s-3)))^{1/2}$.
This bounds the algebra and embedding constants at
orders $s$ and $s-1$, by Fourier convolution and
$\sum_{k\ge1}k^{-2(s-1)}\le1+1/(2s-3)$.
Define the momentum $\Pi=e^{2P}Q_\tau$.
The normalization of the limits uses the same $\tau_0$
for all solutions, so the constants in the expansions
are the intercepts relative to $r$.

\begin{proposition}
\label{prop:low-velocity-difference}
Fix $s\ge3$, $\tau_0$ and $v_*$. Set
$\delta=\min(v_*,1-v_*)/(16A_s)$.
Every smooth periodic vacuum datum with
$\|Z(\tau_0)\|_{H^s}<\delta$ and the Gowdy momentum
constraint evolves to $\tau=\infty$. In a bounded
set of the initial $H^s$ orbit fields, the map from
initial orbit data, velocities and one metric anchor
to
\begin{align*}
 (v,p,q)&=\lim_{r\to\infty}(P_\tau,P-rv,Q),\\
 (\lambda_\infty,\Pi_\infty)&=
                 \lim_{r\to\infty}(\lambda-rv^2,\Pi)
\end{align*}
is locally Lipschitz into
$(H^{s-1})^2\times H^s\times H^{s-1}\times H^{s-2}$.
Take the norm of the input difference to be
\[
 \mathcal I=\|Z(\tau_0)-\widehat Z(\tau_0)\|_{H^s}
      +\|(P,Q,\lambda)(\tau_0)-
                   (\widehat P,\widehat Q,\widehat\lambda)(\tau_0)\|_{H^s}.
\]
One may use the initial $H^{s+1}$ orbit fields, $H^s$
velocities and the scalar metric anchor to bound
$\mathcal I$. Here $v$ is the limiting hyperbolic
speed, with $0<v<1$, and
\begin{equation}
 \partial_\theta\lambda_\infty
 =
 2v\,\partial_\theta p+2\Pi_\infty\,\partial_\theta q .
 \label{eq:tail-momentum}
\end{equation}
\end{proposition}

\begin{proof}
The first-order equations are
\begin{align*}
 V_\tau &= e^{-\tau}W_\theta+X^2-Y^2,\\
 W_\tau &= e^{-\tau}V_\theta-W,\\
 X_\tau &= e^{-\tau}Y_\theta-V X+W Y,\\
 Y_\tau &= e^{-\tau}X_\theta+(V-1)Y-W X .
\end{align*}
Use the squared sum of the component $H^s$ norms
for the energy. Each paired principal spatial
term cancels by periodic integration by parts.
For $D=(W,X,Y)$ and $c=\min(v_*,1-v_*)$,
the diagonal terms give damping at least
$c\|D\|_{H^s}^2$ in the half-energy derivative.
Every nonlinear source contribution to it is
bounded in sum by
$3A_s\|Z\|_{H^s}\|D\|_{H^s}^2$, using the Sobolev
algebra inequality. Smallness therefore gives
\begin{equation}
 \|Z(\tau)\|_{H^s}^2+
             c\int_{\tau_0}^{\tau}\|D(\sigma)\|_{H^s}^2\,d\sigma
     \le \|Z(\tau_0)\|_{H^s}^2.
 \label{eq:low-velocity-energy}
\end{equation}
Sobolev embedding also gives
$0<v_*/2<V< (1+v_*)/2<1$.
The product-source system with symmetric
spatial coupling continues on every finite interval
by the energy estimate. Reconstruct $P$ and $Q$
from $P_\tau=V$ and $Q_\tau=e^{-P}X$,
with their initial values. The definitions of $W$
and $Y$ persist by the $W,Y$ equations and the initial
derivative identities. The usual
propagation of the momentum constraint gives
the periodic metric.

In the damped three-component subsystem,
the $X,Y$ spatial terms still cancel.
At order $s-1$ the remaining spatial forcing
of $W$ is bounded by $e^{-\tau}\|V-v_*\|_{H^s}$.
The nonlinear terms in these three equations
have two damped factors in the half-energy
estimate and a coefficient of size $2A_s\delta$.
Hence for some $\beta>0$ one has
$\|D(\tau_0+r)\|_{H^{s-1}}\le C e^{-\beta r}$.
This also follows by applying the scalar
damped differential inequality to
$\|D\|_{H^{s-1}}$ with input $e^{-\tau}\delta$.

For the difference energy at order $s$, write
$\Delta Z=Z-\widehat Z$. Every nonlinear term
in the half-energy estimate is bounded by
\[
 4A_s\delta\|\Delta D\|_{H^s}^2
  +3A_s(\|D\|_{H^s}+\|\widehat D\|_{H^s})
           \|\Delta Z\|_{H^s}\|\Delta D\|_{H^s}.
\]
The first term is absorbed by damping. Apply
Young's inequality to the second, retaining
damping and a coefficient bounded by
$C(\|D\|_{H^s}^2+\|\widehat D\|_{H^s}^2)$
on the difference energy.
This coefficient is time-integrable by
\eqref{eq:low-velocity-energy}. Consequently
$\sup_r\|\Delta Z\|_{H^s}\le C\mathcal I$.
The equation for $V_\tau$ now gives
$\|\Delta V_\tau\|_{H^{s-1}}
 \le C\mathcal I e^{-\beta r}$, reducing $\beta$
if necessary. The same integrable bound holds for
each solution's $V_\tau$ without $\mathcal I$.
In particular,
\[
 \|v-\widehat v\|_{H^{s-1}}
       +e^{\beta r}\|\Delta V(r)-(v-\widehat v)\|_{H^{s-1}}
 \le C\mathcal I .
\]
The limits of $P-rv$ and of $\lambda-rv^2$ follow
by integrating $V-v$ and
$\lambda_\tau-v^2=V^2-v^2+|D|^2$ respectively.
The differences of these integrands are
$O(\mathcal I e^{-\beta r})$ in $H^{s-1}$.

For the transverse limits, write
$P=P(\tau_0)+v_*r+B$, where
$\|B\|_{H^s}\le\delta r$.
The Sobolev algebra exponential series bounds
$\|e^{\pm B}\|_{H^s}\le C e^{A_s\delta r}$.
The difference of any two such exponentials
has the additional factor
$C(1+r)\mathcal I$, by the product formula
for the difference and
$\|P-\widehat P\|_{H^s}\le(1+Cr)\mathcal I$.
Thus $Q_\tau=e^{-P}X$ and its difference
are integrable in $H^s$ since $A_s\delta<v_*$.
The equation
$\Pi_\tau=e^{-2\tau}\partial_\theta(e^{2P}Q_\theta)$
and the uniform bound for $Q$ in $H^s$ give
integrable bounds for $\Pi_\tau$ and its difference in
$H^{s-2}$, using $A_s\delta<1-v_*$.
The difference estimate has at most an additional
polynomial factor in time. These estimates give the
two remaining limits and their Lipschitz bounds.

Finally $\lambda_\theta=2(V P_\theta+\Pi Q_\theta)$.
Subtract $2rv\,v_\theta$ and take the limit
in $H^{s-2}$ using the exponential remainders.
This gives \eqref{eq:tail-momentum}.
A periodic spatial primitive of the initial
momentum equation bounds the initial $\lambda$
norm by the initial orbit fields, velocities
and its value at a fixed angle.
\end{proof}

\begin{proposition}
\label{prop:from-observations-to-tail}
There are smooth inhomogeneous rank-two
backgrounds in the region of
Proposition~\ref{prop:low-velocity-difference}
for which the finite-snapshot data of
Theorem~\ref{th:fixed-rank-two} determine all five
asymptotic fields with Lipschitz dependence
in the stated norms. One can also use the
two-slice data of Proposition~\ref{prop:two-slice-open}
on an open set of such backgrounds. If neither
initial velocity is supplied, an explicit family
in the same low-velocity region admits recovery
from one full later orbit snapshot and a finite
set of early low-mode orbit coefficients.
\end{proposition}

\begin{proof}
Put $a=e^{-\tau_0}$. Start from
$P=v_*\tau+P_0$, $Q=$ constant.
Take a nearby nonvertical target geodesic
$\gamma$ and replace its scalar wave parameter
$v_*\tau+P_0$ by that parameter plus a small
nonconstant smooth even scalar wave.
The data at $a$ approach the vertical
homogeneous data in the required norms as both
perturbations vanish. Choose the geodesic and
wave so its tangent has nonzero $e_P$ component
on the range at $a$. This is the background
of Proposition~\ref{prop:two-slice-open}.
A small even perturbation of the initial
$Q_t$ transverse to the scalar wave direction
makes the target differential rank two on an open
arc where $P_\theta\ne0$. All orbit fields
and velocities at $a$ remain even, so the
momentum constraint still holds. The finite-time
inverse persists on a neighborhood. The tail
condition also persists by strict smallness.
Composition of the finite-time inverse with
Proposition~\ref{prop:low-velocity-difference}
gives the claim for fixed initial known fields.
Local variation of those fields is allowed as
well, by the upper Cauchy estimate and
the same fixed-background lower bound for
the missing velocity.

For the two-velocity case, at $a$ take
\[
 P=P_h(a)+\epsilon\cos\theta,\quad
 Q=\epsilon\sin\theta,\quad P_t=-v_*/a,\quad Q_t=0,
 \qquad P_h(t)=P_h(a)-v_*\log(t/a).
\]
These data satisfy the periodic momentum constraint.
They have rank two on open arcs and their
$\|Z(\tau_0)\|_{H^s}$ is $O(|\epsilon|)$.
At order one in $\epsilon$ the two displacement
perturbations are $p_1(t)\cos\theta$ and
$q_1(t)\sin\theta$. Their amplitudes start at one
with derivative zero at $a$.
At order two the orbit fields have Fourier modes
only up to order two. These claims follow by
differentiating the equations about the homogeneous
background once and twice.
Choose a sufficiently short rational delay
$d=\pi/L$ with $L>2$ so that $p_1q_1>0$ up to
$c=a+d$. The angle of the rotation in
Proposition~\ref{prop:both-ray-holonomy} has circle
mean $-2\int_a^c\overline{e^P Q_\theta}\,dt$.
In the sum over $L$ shifted rotations all nonzero
Fourier modes of order at most two cancel.
The zero mode of $e^PQ_\theta$ at second order
is $\epsilon^2 e^{P_h}p_1q_1/2$.
Consequently the closed-transport angle equals
\[
 -L\epsilon^2\int_a^c e^{P_h(t)}p_1(t)q_1(t)\,dt
   +O(\epsilon^3)
\]
uniformly on the circle. This is separated from
zero modulo $2\pi$ for small nonzero $\epsilon$.
The rational holonomy criterion and a suitably
early finite set of orbit coefficients give the
stable regular-time inverse for both velocities.
Proposition~\ref{prop:low-velocity-difference}
then gives the same asymptotic conclusion.
Alternatively the same initial curves satisfy the
independence condition of
Theorem~\ref{th:full-metric-gramian}, allowing a
short irrational delay when the terminal metric
gradient is also supplied.
\end{proof}

The estimate applies to the low-velocity region
defined at the recovered initial slice.
It does not assert control through spike
formation or a high-velocity transition.
